\section{METODOLOGIA}

Esta seção apresenta a metodologia usada para o desenvolvimento deste trabalho, isso inclui: a descrição do estudo realizado, o método de construção do trabalho, assim como as métricas utilizadas para a discussão dos resultados dos testes no produto construído.

\subsection{ACERCA DAS CATEGORIAS DA PESQUISA}

As pesquisas podem possuir diversas categorias na metodologia científica, algumas delas são: quanto a finalidade, quanto a objetivos, quanto a procedimentos e quanto a natureza da abordagem. Desse modo, podemos definir exemplos dessas cateogrias relevantes para este texto. Logo, quanto a finalidade há a pesquisa aplicada que possui foco na aplicação, utilização e consequências práticas do conhecimento \cite{gil2008metodos}. Quanto aos objetivos destaca-se a pesquisa descritiva que tem como objetivo primordial a descrição das características de determinada população ou fenômeno \cite{gil2002como}. Quanto aos procedimentos: a pesquisa experimental, determina um objeto de estudo e váriaveis que são capazes de influenciá-lo e define formas de controle e observação dos efeitos que a manipulação dessas variáveis possuem no objeto \cite{gil2008metodos}. E, quanto a natureza da abordagem é de destaque a qualitativa onde a análise é sistemática e compreensiva, usada usualmente em estudos de caso e pesquisa-ação, pois não há fórmula para orientar o pesquisador, assim também precisa ser maleável \cite{gil2008metodos}.

Assim o trabalho se encaixa nas categorias: pesquisa aplicada porque visa a produção de direcionamentos para a implementação das estratégias de sincronização; na descritiva quanto aos seus objetivos de observar os fenômenos (as diferentes estratégias de sincronização em contextos distintos) sem interferir; e, experimental quanto aos procedimentos aplicados, isso é, manipulação dos parâmetros de teste para a simulação de contextos diversos com a coleta de dados qualitativos \cite{gil2008metodos}.


\subsection{ACERDA DO MÉTODO CIENTÍFICO UTILIZADO}

O método científico da pesquisa esclarece os procedimentos lógicos que foram usados na pesquisa \cite{gil2008metodos}, assim, o método escolhido foi o hipotético-dedutivo; este foi escolhido pois define uma abordagem experimental para explicar um fenômeno. Isso se alinha com o objetivo do trabalho de coletar métricas de testes sobre um produto desenvolvido para alcançar um conjunto de postulados que governam os fenômenos analisados na discussão \cite{gil2008metodos}. As métricas utilizadas para comparação entre as diferentes estratégias estão nas categorias: acurácia de sincronização, latência e propagação, resiliência da rede e adaptabilidade, eficiência de transmissão de mensagem e utilização de recursos (memória e processamento). Assim as métricas consistem:

\begin{itemize}[leftmargin=1.5cm]
  \item Porcentagem de consistência de dados: porcentagem de \en{nodes} (ou nós, máquinas na rede que consomem e produzem mensagens) que possuem dados no mesmo estado após um período definido de tempo;
  \item Porcentagem de resolução de conflitos: a porcentagem em que os conflitos são resolvidos corretamente; 
  \item Latência de sincronização: tempo que a mudança em um node leva para alcançar os demais;
  \item Tolerância de \en{peer churn} (conexão e desconexão de nodes na rede): degradação de performance com a conexão e desconexão de nodes frequentes;
  \item Tempo de recuperação: tempo que leva para ressincronizar após falhas;
  \item Excedente de mensagem: razão de mensagens de sincronização para dados atualizados;
  \item Uso de banda larga por \en{node}: quantidade de banda larga usada por por node por sincronização;
  \item Utilização de memória: uso de memória em \en{buffers} e metadados;
  \item Uso de processamento: processamento de mensagens e resolução de conflitos.
\end{itemize}

Essas métricas vão ser acessadas através de testes realizados em uma rede usando libp2p (sistema modular de protocolos, especificações e bibliotecas para desenvolvimento \en{Peer-to-Peer}) para simular um ambiente \en{Peer-to-Peer}. Assim, para o ambiente de simulação de nodes cabe observar os seguintes casos para a coleta de dados: atualização de dados, resolução de conflitos, publicar ou convergir estados (em busca de consistência). A aleatoriedade é importante para testar cenários diversos, logo irá ser implementada nesse ambiente por: variação de atrasos de conexão, conexões/desconexões aleatórias, agrupamento de \en{nodes} em sub redes. Por fim, essas estratégias foram incorporadas na aplicação real.

